\section{PTP Timestamping Integration}
\subsection{Timestamp Authority}
The STM32 nodes are the authoritative timestamp source. Each node disciplines its local clock to the PTP grandmaster using hardware timestamping in the Ethernet peripheral. The \texttt{ts} field of each burst is the time of its first sample, computed in firmware from a hardware timer that shares the PTP clock's 100~MHz source, at the centre of that sample's accumulated ADC burst; it is not broker receive time or Telegraf processing time.

Telegraf's \texttt{json\_time\_key} and \texttt{json\_time\_format = "unix\_ns"} configuration ensures this timestamp is used as the InfluxDB timestamp verbatim.

\subsection{Timescale}
PTP time is TAI. The firmware converts it to UTC by subtracting a fixed 37~s offset before publishing, as it does not yet parse the UTC offset advertised in the grandmaster's Announce messages. A future leap second would therefore require a firmware update. Before the servo first locks, the node's clock counts from boot; such early bursts carry \texttt{sync} = 0.

\subsection{Clock Domain Consistency}
All nodes share a single PTP clock domain. The historian server's system clock must be on UTC, synchronised to the same grandmaster or to a sufficiently accurate NTP source, so that Flux query time ranges expressed in wall-clock time correspond correctly to stored point timestamps.

\subsection{Out-of-Order Writes}
MQTT QoS 0 with no ordering guarantees means Telegraf may deliver points to InfluxDB with timestamps older than recently written points. InfluxDB handles out-of-order writes within a shard natively. Default shard duration in InfluxDB v2 is one week; behaviour at shard boundaries should be verified if a node reconnects after an outage longer than one week, though this is unlikely in a lab testbed context.

\subsection{PTP Status Logging}
Each burst carries a \texttt{sync} field, set to 1 when the node's PTP servo was locked at capture. The servo reports lock after eight consecutive offsets within 1~\textmu s of the grandmaster, with a Sync message received within the last 3~s. Bursts are only started on data whose timestamp is valid (no clock step or rate change during capture), and are restarted if a clock step occurs mid-collection.

\paragraph{Deferred} The original design published PTP diagnostics (\texttt{offset\_ns}, \texttt{path\_delay\_ns}, and \texttt{sync\_state} with LOCKED, ACQUIRING, and HOLDOVER states) to a separate \texttt{ptp\_status} subsystem. This was replaced by the per-burst \texttt{sync} field. The servo offset and path delay are printed on each node's debug UART but not published, so the historian cannot provide post-hoc offset traces.

\subsection{Timestamp Quality Validation}
Before cross-node comparisons are drawn from the historian, \texttt{sync} should be confirmed as 1 for all nodes across the relevant time window. Alignment between nodes can be measured independently of the historian by comparing the 1 pulse-per-second outputs of the nodes on an oscilloscope. This validation is a prerequisite for any result that relies on cross-node temporal alignment.
